% -------------------------------------------------------------------------
% WBS ID: RES-MOD-ASM
% Deliverable: Assumptions, Applicability and Model Limitations
% Status: In progress
% Evidence status: Local 3D baseline assumptions and deferred physics established
% Review status: Not reviewed
% Last updated: 2026-07-18
% Dependencies: RES-MOD-EQS, RES-MOD-SRC, RES-MOD-IBC, RES-PHY-SCL
% Downstream interfaces: RES-NUM-SPEC, RES-VER-MET, RES-STR-FSI
% -------------------------------------------------------------------------

\section{RES-MOD-ASM --- Assumptions, Applicability and Model Limitations}
\label{sec:res-mod-asm}

\subsection{Purpose and Scope}
\label{sec:res-mod-asm-purpose}

% TODO: Define what this Deliverable covers and explicitly excludes.

This Deliverable records the assumptions under which the local
three-dimensional URANS--VOF hydrodynamic model may be used as the
baseline tsunami--barrier interaction model. It also identifies
physical regimes that require deferred extensions before claims can be
made.

\subsection{Research Questions}
\label{sec:res-mod-asm-research-questions}

% TODO: State the questions that must be answered before the Deliverable
% can be accepted.

The local-model applicability questions are:

\begin{enumerate}
  \item Which physical effects are included in the baseline local
    hydrodynamic model?
  \item Which effects are deliberately deferred rather than assumed
    negligible?
  \item Which candidate geometries trigger porous, compressible-air,
    structural-failure, scour or sediment extensions?
  \item Which claims may be made before those extensions and validation
    data exist?
\end{enumerate}

\subsection{Terminology and Definitions}
\label{sec:res-mod-asm-definitions}

% TODO: Define the technical terminology, symbols, quantities and
% classifications used in this Deliverable.

\subsection{Literature Evidence}
\label{sec:res-mod-asm-literature-evidence}

% TODO: Record evidence from peer-reviewed papers, authoritative datasets,
% standards, technical reports and primary documentation.
%
% Do not create a detached literature-summary list. Explain how each source
% supports, contradicts or limits a project decision.

\subsubsection{del Jesus et al. (2012)}
\label{sec:res-mod-asm-evidence-deljesus-2012}

\textcite{deljesus2012porous} demonstrated that unresolved porous
barrier material requires volume averaging and additional
porous-matrix resistance. This formulation is relevant to rubble
mounds, densely packed armour units, aggregate-filled structures, and
other media whose pore geometry is smaller than the resolved
computational scale.

\subsubsection{Istrati and Buckle (2019)}
\label{sec:res-mod-asm-evidence-istrati-2019}

\textcite{istrati2019trappedair} showed that enclosed or poorly
ventilated air volumes can alter local pressure, uplift, and
overturning moment. The study does not imply that compressible flow is
required for every tsunami-impact calculation. It identifies a
geometry-dependent failure regime for the incompressible-air
assumption.

\subsection{Applicability Boundary of the Local Three-Dimensional Model}
\label{sec:res-mod-asm-local-applicability}

The selected incompressible two-phase Navier--Stokes formulation is the
baseline local model for rigid, impermeable, and well-ventilated
barriers. Several alternative or extended formulations were considered
but are not included in the initial model.

\subsubsection{Considered and Deferred Formulations}
\label{sec:res-mod-asm-local-considered-formulations}

\begin{table}[htbp]
  \centering
  \small
  \renewcommand{\arraystretch}{1.25}
  \begin{tabular}{
      p{0.22\textwidth}
      p{0.30\textwidth}
      p{0.37\textwidth}
    }
    \hline
    Formulation
    &
    Reason for consideration
    &
    Project decision
    \\
    \hline
    Compressible two-phase Navier--Stokes
    &
    Represents trapped-gas compression, acoustic effects, and
    compressibility
    &
    Deferred while the design space remains open and ventilated;
    activated only when a trapped-air screening gate fails
    \\
    VARANS
    &
    Represents unresolved internal porous flow through porosity and
    matrix-resistance closures
    &
    Deferred while the initial barrier is impermeable or contains
    explicitly resolved openings
    \\
    Three-dimensional nonhydrostatic multilayer equations
    &
    Provides some vertical resolution at lower dimensional
    complexity
    &
    Not selected owing to limited evidence for local pressure,
    traction, and general moving-boundary FSI
    \\
    Three-dimensional hydrostatic multilayer equations
    &
    Provides vertically layered velocity and density fields
    &
    Rejected for impulsive local impact as the pressure field
    remains hydrostatic
    \\
    Inviscid Euler or potential-flow equations
    &
    Can represent selected nonbreaking wave and body-interaction
    problems
    &
    Rejected for the operational impact model as viscosity,
    turbulent separation, wake formation, and dissipative
    mechanisms are central outputs
    \\
    Molecular-viscosity-only Navier--Stokes
    &
    Retains the parent conservation laws without an explicit
    unresolved-turbulence closure
    &
    Retained only as a diagnostic case; not accepted as the
    high-Reynolds-number operational model
    \\
    SPH or MPS
    &
    Handles highly deformed free surfaces through mesh-free
    numerical formulations
    &
    Deferred to numerical-method selection as these are
    discretisation frameworks rather than alternative conservation
    laws
    \\
    \hline
  \end{tabular}
  \caption{Alternative and extended formulations considered for the
  local three-dimensional model}
  \label{tab:res-mod-asm-local-considered-formulations}
\end{table}

\subsection{Governing Relations and Quantitative Definitions}
\label{sec:res-mod-asm-governing-relations}

% TODO: Record relevant equations, dimensionless groups, empirical models,
% extraction rules or quantitative definitions.
%
% Use align environments for displayed equations.
% Every displayed equation must have a unique \label{eq:...}.
% Do not place punctuation after displayed equations.

\subsection{System Boundary}
\label{sec:res-mod-asm-system-boundary}

% TODO: Complete this RES-MOD Domain-specific subsection.

\subsection{State Variables and Parameters}
\label{sec:res-mod-asm-state-variables-parameters}

% TODO: Complete this RES-MOD Domain-specific subsection.

\subsection{Continuous Governing Model}
\label{sec:res-mod-asm-continuous-governing-model}

% TODO: Complete this RES-MOD Domain-specific subsection.

\subsection{Source Terms and Closures}
\label{sec:res-mod-asm-source-terms-closures}

% TODO: Complete this RES-MOD Domain-specific subsection.

\subsection{Initial, Boundary and Interface Conditions}
\label{sec:res-mod-asm-initial-boundary-interface-conditions}

% TODO: Complete this RES-MOD Domain-specific subsection.

\subsubsection{Initial Local-Model Assumptions}
\label{sec:res-mod-asm-local-initial-assumptions}

The initial local model adopts the following assumptions:

\begin{enumerate}
  \item water and air are Newtonian fluids;
  \item both phases are incompressible;
  \item the phases are immiscible;
  \item heat transfer, evaporation, condensation, and phase change
    are negligible;
  \item the barrier is initially rigid and impermeable;
  \item candidate geometries are sufficiently ventilated to prevent
    material trapped-air compression;
  \item seabed morphology is fixed during one hydrodynamic
    simulation;
  \item debris, sediment transport, scour, cavitation, and structural
    fracture are excluded from the baseline model;
  \item turbulence effects are represented through an attached URANS,
    LES, or hybrid closure;
  \item structural pressure and traction are retained even when the
    structure is fixed.
\end{enumerate}

These assumptions constrain the initial design space rather than assert
that the excluded phenomena are universally negligible.

\subsubsection{Porous-Medium Activation Gate}
\label{sec:res-mod-asm-local-porous-gate}

A porous-region model shall be considered where:

\begin{align}
  \ell_{\mathrm{pore}}
  &\ll
  \Delta_{\mathrm{local}}
  \label{eq:res-mod-asm-porous-scale-separation}
\end{align}

where $\ell_{\mathrm{pore}}$ is a representative pore scale and
$\Delta_{\mathrm{local}}$ is the local resolved mesh scale.

Where individual slots, channels, perforations, or openings are design
variables and can be resolved explicitly, the ordinary Navier--Stokes
equations shall remain active.

Where the internal porous structure cannot be resolved, a local
resistance term may be attached:

\begin{align}
  \mathbf{f}_{\mathrm{por}}
  &=
  -
  \left[
    a\mathbf{u}_{\mathrm{f}}
    +
    b
    \left\lVert
    \mathbf{u}_{\mathrm{f}}
    \right\rVert
    \mathbf{u}_{\mathrm{f}}
    +
    c
    \frac{\partial\mathbf{u}_{\mathrm{f}}}{\partial t}
  \right]
  \label{eq:res-mod-asm-porous-resistance}
\end{align}

The coefficients $a$, $b$, and $c$ shall be calibrated for the
selected material and shall not be transferred directly from unrelated
porous structures.

\subsubsection{Compressible-Air Activation Gate}
\label{sec:res-mod-asm-local-compressible-air-gate}

A candidate geometry shall be flagged for compressible-air assessment
where:

\begin{enumerate}
  \item an initially connected air region becomes isolated during
    impact;
  \item the trapped volume persists for a material fraction of the
    loading interval;
  \item water intrusion causes a material reduction in gas volume;
  \item the venting timescale is comparable with or greater than the
    impact timescale;
  \item cavity pressure contributes materially to the resultant force
    or overturning moment.
\end{enumerate}

A preliminary screening estimate may use:

\begin{align}
  p_{\mathrm{a}}
  V_{\mathrm{a}}^{\,n}
  &=
  p_{\mathrm{a},0}
  V_{\mathrm{a},0}^{\,n}
  \label{eq:res-mod-asm-trapped-air-polytropic-screening}
\end{align}

where $n$ lies between the isothermal and adiabatic limits.

Where the predicted gas-pressure increment is comparable with the local
hydrodynamic pressure, the geometry shall be ventilated, excluded from
the baseline design space, or assessed using a compressible-gas
extension.

\subsubsection{Current Validity Statement}
\label{sec:res-mod-asm-local-validity-statement}

The selected local equation set is applicable to:

\begin{itemize}
  \item rigid or weakly deforming barriers;
  \item impermeable barriers;
  \item explicitly resolved openings and flow passages;
  \item open or adequately ventilated geometries;
  \item water--air free-surface impact without material gas
    compression;
  \item fixed-bed hydrodynamic evaluation;
  \item one-way structural load transfer;
  \item two-way FSI where the structural model and moving-interface
    treatment are activated.
\end{itemize}

The baseline formulation shall not be claimed to reproduce unresolved
porous flow, trapped-air compression, debris impact, movable sediment,
scour, cavitation, fracture, or progressive structural failure without
the corresponding additional physical model.

\subsection{Model Coupling Requirements}
\label{sec:res-mod-asm-model-coupling-requirements}

% TODO: Complete this RES-MOD Domain-specific subsection.

\subsection{Sufficient Conditions for Validity}
\label{sec:res-mod-asm-sufficient-conditions-validity}

% TODO: Complete this RES-MOD Domain-specific subsection.

\subsection{Candidate Approaches}
\label{sec:res-mod-asm-candidate-approaches}

% TODO: Define the credible candidate approaches considered.

\subsection{Critical Comparison}
\label{sec:res-mod-asm-critical-comparison}

% TODO: Compare candidates using physically and project-relevant criteria.
% Do not select an approach solely on popularity or implementation ease.

\subsection{Assumptions and Applicability}
\label{sec:res-mod-asm-assumptions}

% TODO: State assumptions and the conditions under which they are valid.

The local model assumes high-Reynolds-number water--air impact flow
around a fixed or weakly moving defence geometry. The operational
closure is URANS $k$--$\omega$ SST; LES is retained for selected
high-fidelity comparison cases. Source roles: `3D-SST`,
`3D-RANSVOF`, `3D-DEFERRED`. Supporting sources:
\textcite{menter1994twoequation}, \textcite{kozelkov2022turbulence},
\textcite{deljesus2012porous} and \textcite{istrati2019trappedair}.
Project-specific conclusion: turbulence, free-surface deformation and
hydrodynamic loading are baseline phenomena, while unresolved porous
media and trapped-air compressibility are activation-gated extensions.

\subsection{Limitations and Failure Conditions}
\label{sec:res-mod-asm-limitations}

% TODO: State where the evidence, model or selected approach ceases to be
% reliable or applicable.

The baseline model ceases to be sufficient when a quantity of interest
depends materially on:

\begin{itemize}
  \item unresolved flow through porous material;
  \item isolated trapped-air pockets with material compression;
  \item sediment transport, scour or bed morphology change;
  \item debris impact or floating-body collision;
  \item structural fracture, progressive damage or detached debris;
  \item cavitation or compressible acoustic effects;
  \item certification-level design factors not covered by the
    validation evidence.
\end{itemize}

In those regimes the result may be used only as a hydrodynamic
screening calculation unless the missing physical model is added and
validated.

\subsection{Project-Specific Conclusions}
\label{sec:res-mod-asm-conclusions}

% TODO: Convert the evidence into conclusions specific to the Studentship
% project. Do not merely restate literature findings.

The local 3D model is valid as a baseline high-Reynolds-number
hydrodynamic tsunami-barrier interaction model for run-up, overtopping,
free-surface deformation, pressure, force, moment, and comparative
barrier-performance analysis. It is not valid yet for structural
failure, scour, deformable barriers, detailed aeration, sediment
transport, or final design certification.

\subsection{Selected Approach and Rationale}
\label{sec:res-mod-asm-selected-approach}

% TODO: Record the selected approach only after sufficient evidence exists.
% State the alternatives rejected and the reasons for rejection.

The selected baseline assumption set is:

\begin{quote}
  incompressible immiscible water--air flow, VOF free surface,
  high-Reynolds-number URANS $k$--$\omega$ SST turbulence, fixed
  terrain, fixed rigid impermeable and ventilated barrier geometry, and
  hydrodynamic load extraction without baseline structural failure,
  scour, sediment or trapped-air compression
\end{quote}

\subsection{Unresolved Decisions}
\label{sec:res-mod-asm-unresolved-decisions}

% TODO: Record unresolved questions, required evidence and decision owners.

Unresolved applicability details are the quantitative trapped-air
screening threshold, rough-wall turbulence parameters, porous-extension
activation threshold, deformation threshold for activating FSI, and
case-study-specific validation tolerance for treating the barrier as
rigid.

\subsection{Requirements and Handoffs}
\label{sec:res-mod-asm-requirements}

% TODO: State requirements passed to other Research Domains, Software,
% verification activities or later execution work.

`RES-GEO` shall flag candidate geometries that violate the baseline
ventilation, porosity or rigidity assumptions. `RES-VER-MET` shall
define the validation checks that decide whether the hydrodynamic
baseline is acceptable for candidate screening. `RES-STR-FSI` shall own
any structural-feedback or failure extension.

\subsection{Deliverable Completion Record}
\label{sec:res-mod-asm-completion-record}

% TODO: Record whether the Deliverable is:
%
% - Not started
% - In progress
% - Draft complete
% - Reviewed
% - Accepted
%
% Acceptance requires:
%
% - the research questions to be answered;
% - claims to be supported by appropriate evidence;
% - assumptions and limitations to be explicit;
% - the project-specific conclusion to be stated;
% - downstream requirements to be traceable;
% - unresolved decisions to be closed or formally deferred.
